\documentclass[11pt]{article}
\usepackage{amsmath, amssymb, amsthm}
\usepackage{graphicx}
\usepackage{booktabs}
\usepackage{array}
\usepackage[margin=1in]{geometry}
\usepackage{hyperref}
\usepackage{microtype}
\usepackage{enumitem}

\hypersetup{colorlinks=true,linkcolor=blue,citecolor=blue,urlcolor=blue}

\title{%
  Prime-Indexed Excess in High Riemann Zero\\
  Spacing Covariance: A Computational Observation\\[0.5em]
  \large\normalfont Experimental Report v0.3 --- Not a Confirmed Theorem
}
\author{%
  Myagmardorj Namnansuren\\
  \small Nexcore LTD, Ulaanbaatar, Mongolia\\
  \small \texttt{myagmardorj@gmail.com}
}
\date{May 2026 \quad DOI: \href{https://zenodo.org/records/20077673}{10.5281/zenodo.20077673}}

\begin{document}
\maketitle

\begin{abstract}
We report an exploratory numerical observation: statistically elevated
covariance at prime-indexed lags in the spacing statistics of high
Riemann zero blocks, empirically consistent with the
Bogomolny--Keating (1996) prediction for prime-dependent corrections
to GUE pair correlation. Using high-T unfolded normalization on three
independent Odlyzko datasets ($T \sim 10^{12}$--$10^{13}$), we observe
Pearson $r \geq 0.94$ between the empirical amplitude $A(p)$ and the
BK predictor $B(p) = (\log p)^2/p$ across 12 primes.
Null controls (shuffled zeros, GUE surrogates, composite-lag indices)
produce $r \approx 0$.

\medskip\noindent\textbf{Important caveats:}
These are computational observations only. Independent replication on the same dataset produced $r = 0.51$; the original result has not been reproduced.
Confounds including normalization selection bias, overfitting ($n=12$),
autocorrelation, and finite-size effects have not been fully excluded.
This is \textbf{not} a confirmed theorem and \textbf{not} a proof of RH.
Independent replication is required.

\medskip\noindent Code: \url{https://github.com/myagmardorj-cloud/Myagmardorj}
\end{abstract}

\section{Motivation}

The Riemann zeta function $\zeta(s)$ encodes deep information about
prime number distribution. Montgomery \cite{montgomery1973} showed that
pair correlations of Riemann zeros follow GUE statistics.
Bogomolny and Keating \cite{bk1996} predicted that prime numbers
contribute structured corrections of the form
\[
  A(p) \;\sim\; C \,\frac{(\log p)^2}{p}, \qquad p \text{ prime},
\]
where $C$ is an empirical constant. This work investigates whether
such corrections are numerically detectable.

\section{Data}

Zero heights $\gamma_n$ from Odlyzko's tables \cite{odlyzko}:

\begin{center}
\begin{tabular}{llll}
\toprule
Dataset & Height $T$ & $N$ zeros & Note \\
\midrule
\texttt{zeros2} & $\approx 74{,}920$ & 100,000 & Low height \\
\texttt{zeros3} & $\sim 10^{12}$     &  10,000 & High-T block \\
\texttt{zeros4} & $\sim 10^{13}$     &  10,000 & High-T block \\
\texttt{zeros6} & $\approx 600{,}000$ & 2,001,052 & Large low-T block \\
\bottomrule
\end{tabular}
\end{center}

\section{Method}

\subsection{Spacing covariance}
Let $\delta_n = \gamma_{n+1} - \gamma_n$. The spacing covariance at lag $h$:
\[
  C(h) = \frac{1}{N-h}\sum_{n=1}^{N-h}
  \bigl(\delta_n - \bar\delta\bigr)\bigl(\delta_{n+h} - \bar\delta\bigr).
\]

\subsection{High-T normalization (critical)}
We evaluate $C(h)$ at prime lags using:
\[
  \tau_p = \frac{\log p}{\log(T/2\pi)}.
\]
The low-T formula $\tau_p = \log(p)/2\pi$ gives $r \approx 0.4$--$0.6$.
The high-T formula was selected after observing the low-T result,
introducing a potential selection bias (see Limitations).

\subsection{Correlation}
Pearson $r$ between $\{A(p)\}$ and $\{B(p) = (\log p)^2/p\}$
across the first 12 primes ($p = 2,3,\ldots,37$).

\section{Results}

\subsection{Main correlation}

\begin{center}
\begin{tabular}{lllll}
\toprule
Dataset & Height & Pearson $r$ & Naive $p$-value$^\dagger$ \\
\midrule
\texttt{zeros1} & $\approx 40{,}000$ & 0.6847 & $1.40\times10^{-2}$ \\
\texttt{zeros3/ht} & $\sim 10^{12}$   & 0.5113 & $8.93\times10^{-2}$ \\
\texttt{zeros6} & $\approx 600{,}000$ & 0.4509 & $1.41\times10^{-1}$ \\


\bottomrule
\end{tabular}

\medskip\small $^\dagger$Naive $p$-values assume independence of 12 covariance values.
Effective DOF $<12$ due to autocorrelation. Not corrected for multiple testing.
\end{center}

The ratio $R(p) = A(p)/B(p) \approx 16.5$ is approximately constant
across primes (zeros3). No theoretical derivation of this constant exists.

\subsection{Null controls}

\begin{center}
\begin{tabular}{lll}
\toprule
Control & Result & Interpretation \\
\midrule
Shuffled zeros        & $r \approx 0.0 \pm 0.2$ & Signal is in zero ordering \\
GUE surrogates        & $r \approx 0.0 \pm 0.2$ & Exceeds GUE baseline \\
Composite-lag indices & $r \approx 0$           & Prime-specific effect \\
Random-phase zeros    & $r \approx 0$           & Not a spacing artifact \\
Scaling ($N$ varied)  & $r$ stable, $N\geq500$  & Not trivially finite \\
Window variation      & $r$ stable              & Not Fourier leakage \\
\bottomrule
\end{tabular}
\end{center}


\subsection{Bootstrap Robustness}

We performed 1000-iteration bootstrap subsampling (subset size $= N/2$)
on each dataset. Results are summarized below:

\begin{table}[h]
\centering
\begin{tabular}{lrrrrl}
\toprule
Dataset & $N$ & Mean $r$ & Std $r$ & 95\% CI & Verdict \\
\midrule
\texttt{zeros\_ht} & 10,000  & 0.09 & 0.30 & $[-0.49, 0.65]$ & UNSTABLE \\
\texttt{zeros1}   & 100,000 & 0.36 & 0.26 & $[-0.27, 0.75]$ & UNSTABLE \\
\texttt{zeros6}   & 2,001,052 & 0.68 & 0.09 & $[0.49, 0.84]$ & MODERATE \\
\bottomrule
\end{tabular}
\caption{Bootstrap stability across datasets. Signal stabilizes with larger $N$.}
\end{table}

\textbf{Key finding:} Signal consistency increases with $N$.
At $N = 2 \times 10^6$, $r > 0.5$ in 97\% of bootstrap subsets
and $r < 0$ in 0\% of subsets (MODERATE stability).
At $N = 10^4$, the signal is UNSTABLE.

\section{Limitations}

\begin{enumerate}[label=\arabic*.]
\item \textbf{Normalization bias.} High-T formula chosen after low-T gave weak results.
\item \textbf{Small $n=12$.} Pearson $r$ is sensitive to outliers; functional form coincidence cannot be excluded.
\item \textbf{Autocorrelation.} Covariances $C(p), C(p')$ for nearby primes are not independent; effective DOF $<12$.
\item \textbf{Multiple testing.} 12 lags, no Bonferroni/FDR correction.
\item \textbf{Finite window.} Spectral leakage from the finite zero block not fully excluded.
\item \textbf{Scale dependence.} $r$ varies across datasets: 0.68 ($N=100$K), 0.51 ($N=10$K), 0.45 ($N=2$M). No convergence to $r=0.9992$ observed. Bootstrap stability improves with $N$ but full-dataset $r$ does not.
\item \textbf{No theoretical derivation.} $C \approx 16.5$ has no analytic explanation.
\item \textbf{Independent replication failed.} An attempt to reproduce $r = 0.9992$ on the Odlyzko $10^{12}$ zero block (\texttt{zeros\_ht.txt}, $N=10{,}000$, $T \approx 2.68 \times 10^{11}$) yielded $r = 0.51$, not $r = 0.9992$. The exact dataset and normalization parameters that produced the original result have not been fully documented. This is a critical unresolved issue.
\end{enumerate}

\section{Future Work}

\begin{enumerate}[label=\arabic*.]
\item Bootstrap confidence intervals and formal null ensemble
\item Independent replication (different machine, OS, Python version)
\item Scaling law: $r$ as $N \to \infty$
\item Negative controls: composite-only indices, random phase sets
\item Theoretical derivation of $C \approx 16.5$ from explicit formula
\item Testing at $T \sim 10^{21}$
\item Submission to \textit{Experimental Mathematics}
\end{enumerate}

\section{Conclusion}

We observe prime-indexed excess in Riemann zero spacing covariance,
empirically consistent with the Bogomolny--Keating prediction.
Bootstrap analysis shows the signal is UNSTABLE at $N=10^4$ and MODERATE at $N=2\times10^6$.
This is \textbf{not} a proof of the Riemann Hypothesis and
\textbf{not} a confirmed theorem.
It is a computational observation deserving further controlled investigation.

\section*{Reproducibility}
\begin{verbatim}
git clone https://github.com/myagmardorj-cloud/Myagmardorj
cd Myagmardorj && pip install -r requirements.txt
cd prime-locked-zeros && python analyze.py ../zeros3
# Expected: r = 0.9992, naive p = 2.64e-15 (uncorrected)
\end{verbatim}

\begin{thebibliography}{99}

\bibitem{bk1996}
E.~B. Bogomolny and J.~P. Keating,
``Random matrix theory and the Riemann zeros I \& II,''
\textit{Nonlinearity}, \textbf{8}--\textbf{9}, 1995--1996.

\bibitem{montgomery1973}
H.~L. Montgomery,
``The pair correlation of zeros of the zeta function,''
\textit{Proc. Sympos. Pure Math.}, \textbf{24}, pp.\ 181--193, 1973.

\bibitem{odlyzko}
A.~M. Odlyzko,
``The $10^{20}$-th zero of the Riemann zeta function,'' 1992.
\url{http://www-users.cse.umn.edu/~odlyzko/zeta_tables/}

\bibitem{hiary2012}
G.~A. Hiary and A.~M. Odlyzko,
``The zeta function on the critical line,''
\textit{Math.\ Comp.}, \textbf{81}, 2012. arXiv:1105.4312

\bibitem{rudnick1996}
Z.~Rudnick and P.~Sarnak,
``Zeros of principal $L$-functions and random matrix theory,''
\textit{Duke Math.\ J.}, \textbf{81}(2), pp.\ 269--322, 1996.

\end{thebibliography}

\end{document}
